\section{Engine specification}
\label{sec:qs-engines}

The machine-readable specification is \path{step2/composite_qsk/qsk/theory/engine_spec.json}; this section states its content. Both routes compute the same objects, the two-body kernels $\Upsilon_x$ of Definition \ref{def:qs-defect} with their pole and finite parts and the one-body inputs, by independent integration methods and independently written numerator algebra.

\textbf{Route A} (exact simplex integrals, extending \path{step2/twoletter/engine_a} and \path{step2/composite_qsk/composite/engine_a}): chain numerators from the vertex functions of $I$ of the two-letter engine A and the terms of $S^{(1)}(x)$ of Table \ref{tab:qs-terms} realised as in \path{superconformal/model.py}; the special edge with Schwinger time in $[0,\varepsilon]$ and numerator $-|l+q_\sigma|^2$; the Gaussian loop integral and Theorem \ref{thm:qs-simplex} with exact Dirichlet--log integrals over the ordered subsimplices of the region where $\alpha_\sigma$ is smallest; letter derivatives by Proposition \ref{prop:qs-genfun}; the injected momentum $r$ and its derivatives along $\eta$ and $K_b$; colour stripping and the planar adjacent reading of Remark \ref{rem:qs-colour}. Values in $\mathbb Q[i]$ combinations of $1$, $\log2$, $\log3$ and, for $n\ge3$ vertices, the basis of the Part-C engines, confirmed by a second exact evaluation.

\textbf{Route B} (momentum space, extending \path{step2/twoletter/engine_b} and \path{step2/composite_qsk/composite/engine_b}): explicit momentum-space numerator algebra from an independently written realisation of $S^{(1)}(x)$ (the Clifford data of the two-letter engine B with the field transformations of Chapter 2 written out anew), the special edge as the explicit factor $e^{-\varepsilon|p|^2}-1$, Feynman parameters for the other edges, numerical evaluation at several $\varepsilon$ with $\ge30$ digits and Richardson extrapolation to $\varepsilon^0$, the poles fitted and compared; along $K_b$ at the finite endpoints $T|p|^2\in\{1,4,16\}$; and, on the scalar bilinears and the Konishi descendant, the second evaluation of Rem c9-twoeval: $Q$ applied to the arity-one part of $O[\infty]$ along $\xi$ (Lemma c7-image), the graphs with $S^{(1)}(\xi)$ in the flow on a leg, on the cycle and on a one-loop subgraph of a leg, whose sum must be local and equal to \eqref{eq:qs-decomp}.

\textbf{Tasks.} E0 conventions (the commutators of Proposition \ref{prop:qs-comm} on plane waves, the arrows of Lemma \ref{lem:qs-arrows}, the sign rule of Definition \ref{def:qs-defect} (ii) fixed against G-cubic and G-pole, the canonical form with $\gamma^{\mu\nu}\partial_\mu\partial_\nu\delta$, the twisted coefficient $-\tfrac12$); TQ0 the quadratic free directions of Proposition \ref{prop:qs-freexi} (i) and their analogues along $\eta$ and $K_b$; TQ1 the action side along $\xi$: the quadratic part of $\mathrm{Ob}^{(1)}(\xi)$ in the sectors $(\lambda,\AV)$, $(\lambda,\Ap)$, $(c,\lambda^*)$, the bracket $\{S^{(1)}(\xi),\hbar F_1^{(0)}\}$ with the finite counterterm of Thm c4-finite and its $\beta$-family, and the solution $\hbar F_1^{(1)}(\xi)$ of \eqref{eq:qs-arity1xi} with its family $\theta^\xi_V$, $\theta^\xi_\perp$; TQ2 the kernels $\Upsilon_\xi$ on the ten pairs of field and ghost letters, $n\le3$ derivatives (route B $n\le2$), all leg sets of Table \ref{tab:qs-fieldpairs} and of the census file, with the finite special vertices; TQ3 the pairs of Table \ref{tab:qs-twistcensus} at $\xi=u\otimes e_4$ and the decomposition of Theorem \ref{thm:qs-twistcheck} (c), including the arity-two kernel $\Upsilon_{\xi\xi}$ of the entry $bc$; TQ4 the scalar bilinears along $\xi$: $c_Q$ and the three coefficients of the trace part (Cor c7-Q), with their scheme families; TQ5 along $\eta$: the position identity and $\Upsilon^0_\eta$ on $(\Ap,\Ap)$, $(\Ap,\lambda_\pm)$, and $S_1$ on the bilinears and their weight-$\tfrac52$ descendants; TQ6 along $K_b$: $\Upsilon^0_{K_b}$ for $n\le2$ and $K_1$ on the bilinears and the letters by \eqref{eq:qs-K1dec}; TQ7 the regression along $\Delta$ and $P_b$ by the route of the special edges; TQ8 (after the two-point functions) the inner-product tests of Remark \ref{rem:qs-gram} and Remark \ref{rem:qs-KP}.

\textbf{Gates.} G-E0 the commutators; G-TC0 the free directions; G-cubic Example \ref{ex:qs-cubic}; G-noLog no $\log\varepsilon$ in any defect cycle; G-pole \eqref{eq:qs-pole} and G-log \eqref{eq:qs-logcov} pair by pair at $n\le1$; G-pos the position identities of Proposition \ref{prop:qs-posdec}; G-reg Cor \ref{cor:qs-DP} and the first row of the table of Thm c9-D1 ($4\,\mathrm{lc}-\tfrac43O-\tfrac13\{S^{(1)}(\Delta),O\}$ on the bilinears); G-plane Proposition \ref{prop:qs-plane}; G-R1 $sQ_1(O)=0$ on the bilinears; G-R2, G-R3 (the Konishi relation through $D_1$), G-R5, G-R6 on the bilinears on shell (Lemma \ref{lem:qs-onshell}); G-R4 the twisted closure on the pairs of Table \ref{tab:qs-twistcensus}; G-KT the independence of $T$ of $\Upsilon_{K_b}$ (route B at three endpoints, route A by construction) and the vanishing of the $\log T$ coefficient of the defect graphs as $\varepsilon\to0$; G-KD $K_1\cong-(b,x)D_1$ on the bilinears; G-twist Theorem \ref{thm:qs-twistcheck} (a)--(d); G-Ob the quadratic $\mathrm{Ob}^{(1)}(\xi)$ lies in the image of \eqref{eq:qs-arity1xi}; G-AB agreement of the routes entry by entry; G-colour the planar rule in the realisations $\mathfrak{su}(2)$, $\mathfrak{su}(3)$, $\mathfrak{gl}(3)$.
